\section{Python Scripts}

This annex contains the complete Python scripts utilized via the Mininet API to generate the virtual smart home topology and execute the experimental scenarios.

\subsection{Scenario 1: Operatinal Baseline}
\label{annex:scenario1_code}

\begin{lstlisting}
	"""
	=============================================================================
	Network Forensics Project — Scenario 1: The Operational Baseline
	=============================================================================
	Topology
	--------
	[camera]  ──┐
	[thermostat] ──── [gateway (OVS)] ── (internet / external)
	[attacker]──┘
	
	Behaviour
	---------
	• camera      : Continuous UDP stream → gateway:5005
	(simulates an IP camera video feed, ~500 KB/s)
	NTP sync every 600 s (10 min) → 91.189.91.157:123
	• thermostat  : TCP/ping keepalive every 10 s → gateway
	(simulates an MQTT keep-alive heartbeat)
	TCP temperature report every 360 s (6 min) → gateway
	(simulates an MQTT PUBLISH with a temperature reading)
	NTP sync every 600 s (10 min) → 91.189.91.157:123
	• attacker    : Completely silent (no traffic generated)
	
	Capture
	-------
	tcpdump is launched automatically on ALL interfaces (-i any).
	The pcap file will be written to /tmp/scenario1_baseline.pcap inside
	the Mininet VM / Linux host.
	
	Usage
	-----
	sudo python3 scenario1_baseline.py [--duration SECONDS] [--pcap PATH]
	
	Requirements
	------------
	• Mininet   (http://mininet.org)
	• Open vSwitch (OVS)
	• Python 3.6+
	• iperf3    (for UDP stream)  — install with: sudo apt install iperf3
	=============================================================================
	"""
	
	import argparse
	import sys
	import time
	import subprocess
	
	from mininet.net    import Mininet
	from mininet.node   import OVSSwitch, Controller, OVSController
	from mininet.link   import TCLink
	from mininet.log    import setLogLevel, info, error
	from mininet.cli    import CLI
	from mininet.clean  import cleanup
	
	
	# ---------------------------------------------------------------------------
	# Constants
	# ---------------------------------------------------------------------------
	DEFAULT_DURATION   = 3600         # Seconds the scenario runs before auto-stop (60 min). Default, the user can change it using --duration
	DEFAULT_PCAP_PATH  = "/tmp/scenario1_baseline.pcap"
	CAMERA_PORT        = 5005         # UDP destination port (simulated video feed)
	MQTT_PORT          = 1883         # TCP port (simulated MQTT keepalive target)
	NTP_PORT           = 123          # UDP port for NTP time synchronization
	DNS_PORT           = 53           # UDP port for DNS resolution
	TLS_PORT           = 443          # TCP port for TLS cloud telemetry
	KEEPALIVE_INTERVAL = 10           # Seconds between thermostat keepalives
	TEMPERATURE_INTERVAL = 360        # Seconds between thermostat temperature reports
	NTP_INTERVAL       = 600          # Seconds between NTP sync requests (10 min)
	DNS_INTERVAL       = 60           # Seconds between camera DNS queries (1 min)
	TLS_INTERVAL       = 300          # Seconds between camera TLS telemetry reports (5 min)
	NTP_SERVER         = "91.189.91.157" # Simulated external NTP server IP (ntp.ubuntu.com)
	TLS_CLOUD_SERVER   = "93.184.216.34"  # Simulated api.smartcamera.com (uses example.com IP)
	UDP_BANDWIDTH      = "500K"       # iperf3 UDP bandwidth for camera stream
	
	
	# ---------------------------------------------------------------------------
	# Topology builder
	# ---------------------------------------------------------------------------
	
	def build_topology(pcap_path: str):
	"""
	Create and return a running Mininet network with the smart-home topology.
	
	Returns
	-------
	net : Mininet
	The running network object.
	gateway : OVSSwitch
	The central OVS gateway switch.
	camera, thermostat, attacker : Host
	The three IoT / attacker hosts.
	"""
	info("*** Cleaning up any previous Mininet state\n")
	cleanup()
	
	info("*** Creating Mininet network\n")
	net = Mininet(
	controller=None,
	switch=OVSSwitch,
	link=TCLink,
	autoSetMacs=True,
	autoStaticArp=True,
	)
	
	# -- Switch (gateway) -----------------------------------------------------
	info("*** Adding OVS gateway switch\n")
	gateway = net.addSwitch("s1", cls=OVSSwitch, failMode="standalone")
	
	# -- Hosts ----------------------------------------------------------------
	info("*** Adding hosts\n")
	camera     = net.addHost("camera",     ip="10.0.0.1/24", mac="00:00:00:00:00:01")
	thermostat = net.addHost("thermostat", ip="10.0.0.2/24", mac="00:00:00:00:00:02")
	attacker   = net.addHost("attacker",   ip="10.0.0.3/24", mac="00:00:00:00:00:03")
	
	# -- Links ----------------------------------------------------------------
	info("*** Creating links\n")
	net.addLink(camera,     gateway)   # s1-eth1
	net.addLink(thermostat, gateway)   # s1-eth2
	net.addLink(attacker,   gateway)   # s1-eth3
	
	# -- Start network --------------------------------------------------------
	info("*** Starting network\n")
	net.start()
	
	# Verify basic connectivity (ping all-pairs once)
	info("*** Testing connectivity\n")
	net.pingAll()
	
	return net, gateway, camera, thermostat, attacker
	
	
	# ---------------------------------------------------------------------------
	# Traffic generators
	# ---------------------------------------------------------------------------
	
	def start_camera_udp_stream(camera, gateway_ip: str, duration: int):
	"""
	Camera → gateway : continuous UDP stream using iperf3.
	Simulates a low-bitrate IP camera video feed.
	
	The gateway acts as the iperf3 server; the camera is the client.
	We start the server on the gateway IP (but since the gateway is a switch
	without an IP in this topology, we use the thermostat as the reflector /
	sink instead — a realistic IoT NVR scenario).
	"""
	info("*** [camera] Starting UDP video-feed stream\n")
	
	# iperf3 server on thermostat (acts as NVR / cloud endpoint)
	# We launch it in the background with & so it doesn't block.
	camera.cmd(
	f"iperf3 -s -u -p {CAMERA_PORT} -1 --daemon "
	f"--logfile /tmp/iperf3_server_camera.log"
	)
	
	# Small pause to let the server start
	time.sleep(1)
	
	# iperf3 UDP client — camera streams to thermostat IP as a stand-in NVR
	# (In a real deployment this would be an external NVR/cloud server.)
	thermostat_ip = "10.0.0.2"
	camera.cmd(
	f"iperf3 -c {thermostat_ip} -u -p {CAMERA_PORT} "
	f"-b {UDP_BANDWIDTH} -t {duration} "
	f"--logfile /tmp/iperf3_client_camera.log &"
	)
	info(
	f"    camera → thermostat ({thermostat_ip}:{CAMERA_PORT}) "
	f"UDP {UDP_BANDWIDTH} for {duration}s\n"
	)
	
	
	def start_thermostat_keepalive(thermostat, gateway_ip: str, duration: int):
	"""
	Thermostat → gateway : a lightweight TCP SYN / ping every KEEPALIVE_INTERVAL s.
	Simulates an MQTT PINGREQ keep-alive packet.
	
	The TCP connection is established every 10 seconds from the thermostat to the
	gateway's MQTT port (1883). 
	
	We use a small Python one-liner loop running inside the host's shell so
	the timing is handled in the background independently of the main process.
	"""
	info("*** [thermostat] Starting MQTT keep-alive simulation\n")
	
	# We send a TCP connect (3-way handshake) then immediately close it, which
	# is representative of an MQTT PINGREQ/PINGRESP exchange at the packet level.
	# Using /dev/tcp bash built-in avoids needing extra tools.
	keepalive_script = (
	f"while true; do "
	f"  python3 -c \""
	f"import socket, time; "
	f"s = socket.socket(); "
	f"s.settimeout(3); "
	f"s.connect(('{gateway_ip}', {MQTT_PORT})); "
	f"s.send(b'MQTT_PINGREQ'); "
	f"s.close()\" 2>/dev/null || "
	# Fallback: plain ICMP ping if TCP connection is refused
	f"  ping -c 1 -W 2 {gateway_ip} > /dev/null 2>&1; "
	f"  sleep {KEEPALIVE_INTERVAL}; "
	f"done"
	)
	
	thermostat.cmd(f"bash -c '{keepalive_script}' &")
	info(
	f"    thermostat → gateway ({gateway_ip}:{MQTT_PORT}) "
	f"TCP keepalive every {KEEPALIVE_INTERVAL}s\n"
	)
	
	
	def start_thermostat_temperature_report(thermostat, gateway_ip: str, duration: int):
	"""
	Thermostat → gateway : a simulated MQTT PUBLISH every TEMPERATURE_INTERVAL s.
	Simulates a thermostat sending its current temperature.
	
	Every 6 minutes, the thermostat sends a TCP packet to the gateway's MQTT port (1883).
	"""
	info("*** [thermostat] Starting MQTT temperature report simulation\n")
	
	report_script = (
	f"while true; do "
	f"  python3 -c \""
	f"import socket, time; "
	f"s = socket.socket(); "
	f"s.settimeout(3); "
	f"s.connect(('{gateway_ip}', {MQTT_PORT})); "
	f"s.send(b'MQTT_PUBLISH: topic=home/thermostat/temperature payload=22.5C'); "
	f"s.close()\" 2>/dev/null || "
	f"  ping -c 1 -W 2 {gateway_ip} > /dev/null 2>&1; "
	f"  sleep {TEMPERATURE_INTERVAL}; "
	f"done"
	)
	
	thermostat.cmd(f"bash -c '{report_script}' &")
	info(
	f"    thermostat → gateway ({gateway_ip}:{MQTT_PORT}) "
	f"TCP temperature report every {TEMPERATURE_INTERVAL}s\n"
	)
	
	
	def start_ntp_sync(node, ntp_server_ip: str):
	"""
	Simulates periodic NTP requests from an IoT device to an external NTP server.
	Sends a valid 48-byte NTP client request packet over UDP port 123.
	"""
	info(f"*** [{node.name}] Starting simulated NTP sync to {ntp_server_ip}\n")
	
	ntp_script = (
	f"while true; do "
	f"  python3 -c \""
	f"import socket; "
	f"s = socket.socket(socket.AF_INET, socket.SOCK_DGRAM); "
	f"s.sendto(b'\\x1b' + 47 * b'\\x00', ('{ntp_server_ip}', {NTP_PORT}))\" 2>/dev/null; "
	f"  sleep {NTP_INTERVAL}; "
	f"done"
	)
	
	node.cmd(f"bash -c '{ntp_script}' &")
	
	
	def start_camera_dns_queries(camera, gateway_ip: str):
	"""
	Camera → Local Gateway : periodic UDP DNS queries on port 53.
	Simulates the camera resolving cloud hostnames (api.smartcamera.com,
	stream.smartcamera.com, pool.ntp.org) before making outbound connections.
	
	A minimal valid DNS query for 'api.smartcamera.com' is crafted and sent
	as a raw UDP datagram to the gateway IP on port 53.
	"""
	info(f"*** [camera] Starting simulated DNS queries to {gateway_ip}\n")
	
	# Minimal DNS query wire format for 'api.smartcamera.com' (type A, class IN)
	dns_query = (
	"b'\\x00\\x01'"  # Transaction ID: 1
	" + b'\\x01\\x00'"  # Flags: standard recursive query
	" + b'\\x00\\x01'"  # Questions: 1
	" + b'\\x00\\x00' * 3"  # Answer/NS/Additional RRs: 0
	" + b'\\x03api\\x0csmartcamera\\x03com\\x00'"  # QNAME
	" + b'\\x00\\x01'"  # QTYPE: A
	" + b'\\x00\\x01'"  # QCLASS: IN
	)
	
	dns_script = (
	f"while true; do "
	f"  python3 -c \""
	f"import socket; "
	f"s = socket.socket(socket.AF_INET, socket.SOCK_DGRAM); "
	f"s.sendto({dns_query}, ('{gateway_ip}', {DNS_PORT}))\" 2>/dev/null; "
	f"  sleep {DNS_INTERVAL}; "
	f"done"
	)
	
	camera.cmd(f"bash -c '{dns_script}' &")
	info(
	f"    camera → gateway ({gateway_ip}:{DNS_PORT}) "
	f"UDP DNS query every {DNS_INTERVAL}s\n"
	)
	
	
	def start_camera_tls_telemetry(camera, cloud_server_ip: str):
	"""
	Camera → Internet : periodic TCP connection on port 443.
	Simulates the camera sending TLS cloud telemetry to api.smartcamera.com.
	
	A TCP SYN + TLS ClientHello-like handshake initiation is performed.
	The connection is closed immediately after (no full TLS negotiation
	is performed since the server is simulated), which is sufficient to
	produce the correct packet signature in the pcap capture.
	"""
	info(f"*** [camera] Starting simulated TLS telemetry to {cloud_server_ip}\n")
	
	tls_script = (
	f"while true; do "
	f"  python3 -c \""
	f"import socket; "
	f"s = socket.socket(socket.AF_INET, socket.SOCK_STREAM); "
	f"s.settimeout(3); "
	f"s.connect(('{cloud_server_ip}', {TLS_PORT})); "
	f"s.send(b'\\x16\\x03\\x01\\x00\\xf1\\x01\\x00\\x00\\xed\\x03\\x03'); "
	f"s.close()\" 2>/dev/null; "
	f"  sleep {TLS_INTERVAL}; "
	f"done"
	)
	
	camera.cmd(f"bash -c '{tls_script}' &")
	info(
	f"    camera → cloud ({cloud_server_ip}:{TLS_PORT}) "
	f"TCP TLS telemetry every {TLS_INTERVAL}s\n"
	)
	
	
	def start_thermostat_dns_queries(thermostat, gateway_ip: str):
	"""
	Thermostat → Local Gateway : periodic UDP DNS queries on port 53.
	Simulates the thermostat resolving the MQTT broker hostname
	(mqtt.smartthermostat.com) before establishing its connection.
	
	A minimal valid DNS query for 'mqtt.smartthermostat.com' is crafted
	and sent as a raw UDP datagram to the gateway IP on port 53.
	"""
	info(f"*** [thermostat] Starting simulated DNS queries to {gateway_ip}\n")
	
	# Minimal DNS query wire format for 'mqtt.smartthermostat.com' (type A, class IN)
	dns_query = (
	"b'\\x00\\x02'"  # Transaction ID: 2
	" + b'\\x01\\x00'"  # Flags: standard recursive query
	" + b'\\x00\\x01'"  # Questions: 1
	" + b'\\x00\\x00' * 3"  # Answer/NS/Additional RRs: 0
	" + b'\\x04mqtt\\x0fsmartThermostat\\x03com\\x00'"  # QNAME
	" + b'\\x00\\x01'"  # QTYPE: A
	" + b'\\x00\\x01'"  # QCLASS: IN
	)
	
	dns_script = (
	f"while true; do "
	f"  python3 -c \""
	f"import socket; "
	f"s = socket.socket(socket.AF_INET, socket.SOCK_DGRAM); "
	f"s.sendto({dns_query}, ('{gateway_ip}', {DNS_PORT}))\" 2>/dev/null; "
	f"  sleep {DNS_INTERVAL}; "
	f"done"
	)
	
	thermostat.cmd(f"bash -c '{dns_script}' &")
	info(
	f"    thermostat → gateway ({gateway_ip}:{DNS_PORT}) "
	f"UDP DNS query every {DNS_INTERVAL}s\n"
	)
	
	
	def attacker_activity(attacker):
	"""
	Attacker host: No traffic will be generated.
	The interface is intentionally left idle to reflect a real forensic
	baseline where the attacker has not yet begun their campaign.
	"""
	info("*** [attacker] Node is SILENT — no traffic will be generated\n")
	# Deliberately no commands issued to the attacker host.
	
	
	# ---------------------------------------------------------------------------
	# Capture instructions (printed to stdout for the operator)
	# ---------------------------------------------------------------------------
	
	TCPDUMP_BANNER = """
	╔══════════════════════════════════════════════════════════════════════════════╗
	║              tcpdump CAPTURE — Scenario 1 Baseline                         ║
	╠══════════════════════════════════════════════════════════════════════════════╣
	║                                                                              ║
	║  Capturing ALL traffic on every interface (-i any).                         ║
	║                                                                              ║
	║  The capture file will be saved to:                                          ║
	║    {pcap_path}                                                               ║
	║                                                                              ║
	║  To open the capture afterwards:                                             ║
	║    wireshark {pcap_path}                                                     ║
	║    tshark -r {pcap_path} -T fields -e frame.number -e ip.src -e ip.dst      ║
	║            -e ip.proto -e frame.len                                          ║
	╚══════════════════════════════════════════════════════════════════════════════╝
	"""
	
	
	def print_capture_instructions(pcap_path: str):
	print(TCPDUMP_BANNER.format(pcap_path=pcap_path))
	
	
	# ---------------------------------------------------------------------------
	# Automated in-process tcpdump launch
	# ---------------------------------------------------------------------------
	
	def launch_tcpdump(pcap_path: str) -> subprocess.Popen:
	"""
	Launch tcpdump on all interfaces to capture the full baseline.
	Returns the Popen handle so the caller can terminate it later.
	"""
	iface = "any"
	cmd = [
	"tcpdump",
	"-i", iface,
	"-w", pcap_path,
	"--immediate-mode",
	]
	info(f"*** Launching tcpdump on {iface} → {pcap_path}\n")
	try:
	proc = subprocess.Popen(
	cmd,
	stdout=subprocess.DEVNULL,
	stderr=subprocess.PIPE,
	)
	time.sleep(1)  # Give tcpdump a moment to open the capture file
	return proc
	except FileNotFoundError:
	error("tcpdump not found — install with: sudo apt install tcpdump\n")
	return None
	except PermissionError:
	error("tcpdump requires root privileges. Run the script with sudo.\n")
	return None
	
	
	# ---------------------------------------------------------------------------
	# Main
	# ---------------------------------------------------------------------------
	
	def parse_args():
	parser = argparse.ArgumentParser(
	description="Scenario 1 — Smart-home Operational Baseline (Mininet)"
	)
	parser.add_argument(
	"--duration", "-d",
	type=int,
	default=DEFAULT_DURATION,
	help=f"How long (seconds) to run the scenario before stopping. "
	f"Default: {DEFAULT_DURATION}. Use 0 to drop into interactive CLI.",
	)
	parser.add_argument(
	"--pcap", "-p",
	type=str,
	default=DEFAULT_PCAP_PATH,
	help=f"Output path for the baseline pcap file. "
	f"Default: {DEFAULT_PCAP_PATH}",
	)
	return parser.parse_args()
	
	
	def run_scenario(args):
	setLogLevel("info")
	
	info("=" * 70 + "\n")
	info(" Scenario 1 — Smart-Home Operational Baseline\n")
	info("=" * 70 + "\n")
	
	# Print capture instructions for the operator
	print_capture_instructions(args.pcap)
	
	# ── Build topology ────────────────────────────────────────────────────
	net, gateway, camera, thermostat, attacker = build_topology(args.pcap)
	
	# Determine gateway IP.  Because the gateway is an OVS switch it does
	# not have an IP assigned by Mininet.  We use thermostat as the MQTT
	# broker / NVR proxy target (realistic for a home-network scenario).
	# For keep-alive pings we fall back to ICMP → thermostat.
	gateway_ip = "10.0.0.2"   # thermostat acts as broker / NVR endpoint
	
	# ── Launch tcpdump on all interfaces ─────────────────────────────────
	tcpdump_proc = launch_tcpdump(args.pcap)
	
	# ── Start traffic ────────────────────────────────────────────────────
	info("\n*** Starting traffic generators\n")
	
	attacker_activity(attacker)
	
	start_thermostat_keepalive(thermostat, gateway_ip, args.duration)
	
	start_thermostat_temperature_report(thermostat, gateway_ip, args.duration)
	
	start_thermostat_dns_queries(thermostat, gateway_ip)
	
	start_camera_udp_stream(camera, gateway_ip, args.duration)
	
	start_camera_dns_queries(camera, gateway_ip)
	start_camera_tls_telemetry(camera, TLS_CLOUD_SERVER)
	
	start_ntp_sync(camera, NTP_SERVER)
	start_ntp_sync(thermostat, NTP_SERVER)
	
	# ── Run scenario ─────────────────────────────────────────────────────
	if args.duration == 0:
	info("\n*** Dropping into interactive CLI (duration=0)\n")
	info("    Type 'exit' or Ctrl-D to stop the scenario.\n\n")
	CLI(net)
	else:
	info(f"\n*** Scenario running for {args.duration} seconds …\n")
	info("    Press Ctrl-C to stop early.\n\n")
	try:
	for elapsed in range(args.duration):
	time.sleep(1)
	if (elapsed + 1) % 10 == 0:
	info(f"    [{elapsed + 1:4d}/{args.duration}s] "
	f"Scenario 1 running — "
	f"camera streaming, thermostat active, NTP active, "
	f"attacker silent\n")
	except KeyboardInterrupt:
	info("\n*** Interrupted by user\n")
	
	# ── Teardown ─────────────────────────────────────────────────────────
	info("\n*** Stopping scenario\n")
	
	if tcpdump_proc and tcpdump_proc.poll() is None:
	tcpdump_proc.terminate()
	tcpdump_proc.wait()
	info(f"*** tcpdump stopped — capture saved to {args.pcap}\n")
	
	net.stop()
	info("*** Network stopped\n")
	info(f"\n    Baseline capture: {args.pcap}\n")
	info("    Open with: wireshark " + args.pcap + "\n")
	info("=" * 70 + "\n")
	
	
	if __name__ == "__main__":
	args = parse_args()
	
	# Mininet requires root
	import os
	if os.geteuid() != 0:
	print("[ERROR] This script must be run with sudo / as root.")
	sys.exit(1)
	
	run_scenario(args)
\end{lstlisting}

\subsection{Scenario 2: Network Reconnaissance and Brute-Force Authentication}
\label{annex:scenario2_code}

\subsection{Scenario 3: Unauthorized State Change and Data Exfiltration}
\label{annex:scenario3_code}

\subsection{Scenario 4: Botnet Infection (Outbound DDoS)}
\label{annex:scenario4_code}

